\begin{tablalibro}{2}{L{0.31}L{0.69}}
\caption{Glosario del capítulo}\label{tab:glosario-del-capitulo-1}\\
\toprule
\rowcolor{Papel}
\CabeceraTabla{Término} & \CabeceraTabla{Significado} \\
\midrule
\endfirsthead
\SigueTabla{2}
\toprule
\rowcolor{Papel}
\CabeceraTabla{Término} & \CabeceraTabla{Significado} \\
\midrule
\endhead
\midrule
\endfoot
\bottomrule
\multicolumn{2}{@{}p{\linewidth}@{}}{\NotaLibro{Elaboración propia.}}\\
\endlastfoot

Problema computacional & Necesidad expresable mediante datos y procedimientos computacionales. \\
Pensamiento computacional & Razonamiento orientado a formular y organizar soluciones computacionales. \\
Descomposición & División de un problema en partes manejables. \\
Patrón & Regularidad o estructura común entre situaciones. \\
Abstracción & Selección de características relevantes para un objetivo. \\
Generalización & Extensión de una solución a casos semejantes. \\
Entrada & Dato requerido para efectuar un proceso. \\
Proceso & Transformación u operación sobre los datos. \\
Salida & Resultado producido por el procedimiento. \\
Restricción & Condición que limita la solución. \\
\end{tablalibro}
